\chapter{拓扑的应用}

\section{代数学基本定理}

\begin{lemma}[][lem:11.1]
    $h:S^1\to \mathbb{C}-\{0\}$是零伦的当且仅当$h$可以连续延拓到$g: D^2 \to \mathbb{C}-\{0\}$.
\end{lemma}
\begin{proof}
    $(\impliedby)$
    假设 $h$ 可延拓为 $g: D^2 \to \mathbb{C}-\{0\}$，则$h = g|_{\partial D^2}: S^1 \to \mathbb{C}-\{0\}$.\par
    取常值映射$\setjot[-4]\begin{aligned}[t]
        h_0: S^1 &\to \mathbb{C}-\{0\} \\
        z &\mapsto g(0),0 \in D^2
    \end{aligned}$，
    构造 $\setjot[-4]\begin{aligned}[t]
        H: S^1 \times [0, 1] &\to \mathbb{C}-\{0\} \\
        (z, t) &\mapsto g(tz)
    \end{aligned}$.\par
    则$H(z, 1)=g(z)=h(z)$，$H(z, 0)=g(0)=h_0(z)$，故 $h \simeq h_0$，$h$ 是零同伦的.\par
    $(\implies)$
    假设 $h$ 是零同伦的，即存在常值映射 $c: S^1 \to \mathbb{C}-\{0\}$ 及同伦 $H: S^1 \times [0, 1] \longrightarrow \mathbb{C}-\{0\}$ 使得 $H(z, 1) = h(z)$ 且 $H(z, 0) = c$.
    
    由于 $D^2 \cong (S^1 \times [0, 1]) / (S^1 \times \{0\})$.
    定义$\setjot[-2]\begin{aligned}[t]
        g: D^2 &\to \mathbb{C}-\{0\} \\
        x &\mapsto\begin{cases}
            H\left(\frac{x}{|x|}, |x|\right),&x\neq 0 \\
            c,&x=0
        \end{cases}
    \end{aligned}$\par
    由于 $H(z, 0) = c$ 对所有 $z \in S^1$ 成立，故 $g$ 在 $0$ 处良定义且连续.
    当 $x \in S^1$ 时，$t=1$，则 $g(x) = H(x, 1) = h(x)$.
    所以 $g$ 是 $h$ 的连续延拓.
\end{proof}

\begin{theorem}[代数学基本定理][thm:11.1]
    $f(z) = z^n + a_{n-1}z^{n-1} + \dots + a_1 z + a_0 = 0$ 在 $\mathbb{C}$ 中有根，其中 $a_i \in \mathbb{C}, n \geq 1$.
\end{theorem}

\begin{proof}
    (1) 特殊情况：$|a_{n-1}| + \dots + |a_1| + |a_0| < 1$.\par
    假设 $f(z)=0$ 在 $\mathbb{C}$ 中没有根.\par
    令 $D^2 = \{z \in \mathbb{C} \mid |z| \leq 1\}$，$\partial D^2 = S^1$，
    $\setjot[-4]\begin{aligned}[t]
        g: D^2 &\to \mathbb{C}-\{0\} \\
        z &\mapsto f(z) = z^n + a_{n-1}z^{n-1} + \dots + a_1 z + a_0
    \end{aligned}$.\par
    令 $h = g|_{\partial D^2}: S^1 \longrightarrow \mathbb{C}-\{0\}$.

    由于 $h$ 可延拓为 $g$，故 $h$ 是零同伦的.\par
    构造同伦 $\setjot[-4]\begin{aligned}[t]
        H: S^1 \times [0, 1] &\to \mathbb{C}-\{0\} \\
        (z, t) &\mapsto z^n + t(a_{n-1}z^{n-1} + \dots + a_1 z + a_0)
    \end{aligned}$\par
    我们有$\setjot[-4]\begin{aligned}[t]
        |H(z, t)| &= |z^n + t(a_{n-1}z^{n-1} + \dots + a_1 z + a_0)| \\
    &\geq |z^n| - |t(a_{n-1}z^{n-1} + \dots + a_1 z + a_0)| \\
    &= 1 - t|a_{n-1}z^{n-1} + \dots + a_1 z + a_0| \\
    &\geq 1 - t(|a_{n-1}| + \dots + |a_1| + |a_0|) > 0
    \end{aligned}$

    因此 $H$ 是良定义的且连续.

    $H$ 实现 $h \simeq z^n$，其中$\setjot[-4]\begin{aligned}[t]
        S^1 &\to \mathbb{C}-\{0\} \\
        e^{2\pi i t} &\mapsto e^{2\pi i n t}
    \end{aligned}$.\par
    而$z^n$ 在基本群之间诱导了非平凡同态，所以 $h$ 不是零同伦的.矛盾！ \par
    所以 $f(z)=0$ 在 $\mathbb{C}$ 中有根.\par
    (2) 一般情况，$z^n + a_{n-1}z^{n-1} + \dots + a_1 z + a_0 = 0$.\par
    令 $z=cy$，则$(cy)^n + a_{n-1}(cy)^{n-1} + \dots + a_1(cy) + a_0 = 0$.\par
    取 $c \in \mathbb{R}$，$c$ 足够大使得$\left|\frac{a_{n-1}}{c}\right| + \dots + \left|\frac{a_1}{c^{n-1}}\right| + \left|\frac{a_0}{c^n}\right| < 1$.\par
    方程 $y^n + \frac{a_{n-1}}{c}y^{n-1} + \dots + \frac{a_1}{c^{n-1}}y + \frac{a_0}{c^n} = 0$ 属于情况 (1)，它有根 $y_0 \in \mathbb{C}$.\par
    那么 $z^n + a_{n-1}z^{n-1} + \dots + a_1 z + a_0 = 0$ 有根 $z_0 = c y_0 \in \mathbb{C}$.
\end{proof}

\begin{lemma}[][lem:11.2]
    $f:X\to Y$是零伦的当且仅当$f$可以连续扩张到锥空间$CX$上.
\end{lemma}
\begin{proof}
    $(\implies)$
    若 $f$ 是零伦的，则存在常值映射$\setjot[-6]\begin{aligned}[t]
        c: X &\to Y \\
        x &\mapsto y_0
    \end{aligned}$及同伦$H: X \times [0, 1] \to Y$ 使得 $H(x, 0) = f(x)$ 且 $H(x, 1) = y_0$.\par
    锥空间定义为 $CX = (X \times [0, 1]) / (X \times \{1\})$.
    由于 $H$ 将集合 $X \times \{1\}$ 映射为单点 $y_0$，根据商映射泛性质，$H$ 诱导了连续映射 $F: CX \to Y$，满足 $F([x, t]) = H(x, t)$.\par
    限制在底面 $X \cong X \times \{0\}$ 上，$F([x, 0]) = H(x, 0) = f(x)$，即 $F$ 是 $f$ 的扩张.\par
    $(\impliedby)$
    若 $f$ 能扩张为 $F: CX \to Y$，满足 $F([x, 0]) = f(x)$.\par
    令 $q: X \times [0, 1] \to CX$ 为商映射. 定义 $H = F \circ q: X \times [0, 1] \to Y$.\par
    则 $H$ 连续，且 $H(x, 0) = F([x, 0]) = f(x)$，$H(x, 1) = F([x, 1]) = F(*)$ (常值).\par
    故 $H$ 是 $f$ 到常值映射的同伦，$f$ 是零伦的.
\end{proof}


\section{不动点}

\begin{lemma}[][lem:11.3]
    $S^1 = \partial D^2$，$S^1$ 不是 $D^2$ 的收缩核 (retract).
\end{lemma}

\begin{proof}
    假设 $r: D^2 \longrightarrow S^1$ 是一个收缩.\par
    那么 $r \circ i = id_{S^1}: S^1 \longrightarrow S^1$. \par
    取 $x_0 \in S^1$，则有下述交换图成立： \par
    \begin{tikzcd}[labels={font=\normalsize},row sep=large, column sep=large]
    \makebox[0pt][r]{\textcolor{red}{$x_0\in$\,}}S^1 \arrow[rr, "id_{S^1}"] \arrow[dr, "i"'] & & S^1 \\
    & D^2 \arrow[ru, "r"'] &
    \end{tikzcd}
    \hspace{1em}
    \begin{tikzcd}[labels={font=\normalsize},row sep=large, column sep=large]
    % 左上节点：标注放在左侧 (makebox[0pt][r])
    % \quad 用于增加一点间距
    \makebox[0pt][r]{\textcolor{red}{$\mathbb{Z}\cong$\,}} \pi_1(S^1, x_0) 
      \arrow[rr, "id"] 
      \arrow[dr, "i_*"'] 
    & & 
    % 右上节点：标注放在右侧 (makebox[0pt][l])
    \pi_1(S^1, x_0) \makebox[0pt][l]{\,\textcolor{red}{$\cong \mathbb{Z}$}} 
    \\
    & 
    % 下方节点：标注放在右侧
    \pi_1(D^2, x_0) \makebox[0pt][l]{\,\textcolor{red}{$\cong\{1\}$}}
      \arrow[ur, "r_*"'] 
    & 
    \end{tikzcd} \par
    即$r_* \circ i_* = id: \pi_1(S^1, x_0) \longrightarrow \pi_1(S^1, x_0)$.\par
    由于 $\pi_1(D^2, x_0) = \{1\}$，$r_* \circ i_*$ 是平凡同态，矛盾！
\end{proof}

\begin{theorem}[Brouwer不动点定理][thm:11.2]
\index{B@Brouwer不动点定理}
    若 $f: D^2 \longrightarrow D^2$ 是连续的，则 $f$ 有不动点，即 $\exists\,x_0 \in D^2$，使得 $f(x_0) = x_0$.
\end{theorem}

\begin{proof}
    假设 $\forall x \in D^2, f(x) \neq x$.构造一个连续映射 $g: D^2 \longrightarrow S^1$ 如下：

    \begin{figure}[H]
        \centering
        \includegraphics[width=0.4\textwidth]{example-image}
    \end{figure}

    对于任意 $x \in S^1$，$g(x)$为$x$和$f(x)$的连线与$S^1$的交点.\par
    所以 $g$ 是一个从 $D^2$ 到 $S^1$ 的收缩.\par
    矛盾！故 $f: D^2 \longrightarrow D^2$ 必须有不动点.
\end{proof}
使用Brouwer不动点定理可以证明以下定理：
\begin{theorem}[Frobenius定理][thm:11.3]
    假设 $A$ 是一个 $3 \times 3$ 矩阵.
    \[
        \begin{pmatrix}
        a_{11} & a_{12} & a_{13} \\
        a_{21} & a_{22} & a_{23} \\
        a_{31} & a_{32} & a_{33}
        \end{pmatrix}, \quad a_{ij} \in \mathbb{R}, \quad a_{ij} > 0
    \]
    那么 $A$ 有一个特征值 $\lambda \in \mathbb{R}$，且 $\lambda > 0$.
\end{theorem}

\begin{proof}
    定义$\setjot[-2]\begin{aligned}[t]
        \sigma: \mathbb{R}^3 &\to \mathbb{R}^3 \\
        \begin{pmatrix} x_1 \\ x_2 \\ x_3 \end{pmatrix} &\mapsto A \begin{pmatrix} x_1 \\ x_2 \\ x_3 \end{pmatrix}
    \end{aligned} $\par
    令 $\Omega = \{(x_1, x_2, x_3) \in S^2 \mid x_i \geq 0, i=1,2,3\}$.\par
    则$\Omega \cong D^2$，且$\forall x \in \Omega,\sigma(x) \neq 0$，$\frac{\sigma(x)}{\|\sigma(x)\|} \in \Omega$.\par
    \begin{figure}[H]
        \centering
        \includegraphics[width=0.4\textwidth]{example-image}
    \end{figure}
    故映射$\setjot[-4]\begin{aligned}[t]
        f: \Omega &\to \Omega \\
        x &\mapsto \frac{\sigma(x)}{\|\sigma(x)\|}
    \end{aligned}$有一个不动点 $x^0$.\par
    $\frac{\sigma(x^0)}{\|\sigma(x^0)\|} = x^0 \implies \sigma(x^0) = \|\sigma(x^0)\| \cdot x^0$.\par
    所以 $\|\sigma(x^0)\|$ 是 $A$ 的一个特征值，且满足$\|\sigma(x^0)\| \in \mathbb{R}$，$\|\sigma(x^0)\| > 0$.
\end{proof}

